%% ════════════════════════════════════════════════════════════════════════════
\section{AI and Analytics Layer}
\label{sec:ai}
%% ════════════════════════════════════════════════════════════════════════════

\subsection{Overview and Motivation}
\label{subsec:ai_overview}

The AI analytics layer adds two capabilities beyond raw sensor and
attendance data: (1)~at-risk student detection with natural-language
explanations and (2)~per-course attendance forecasting with trend
classification.  Both pipelines are implemented as on-demand background
tasks triggered by professor page visits, rather than scheduled nightly
jobs, ensuring explanations are available immediately on first access.

The choice of a local LLM (phi3-mini via Ollama) over a cloud API rests
on three requirements: student attendance data must remain within the
university network boundary; the system must operate without internet
access; and inference cost must be zero beyond hardware.  phi3-mini
(3.8~B parameters, $\sim$4\,k-token context) provides adequate language
quality for 3--4 sentence analytical summaries while running on the
Raspberry Pi~4B CPU.

\subsection{At-Risk Student Detection Pipeline}
\label{subsec:at_risk_pipeline}

\subsubsection{Trigger and Lock Mechanism}

The at-risk pipeline is triggered when a professor navigates to
\texttt{/at-risk}, which calls \texttt{GET /api/at-risk}.  The backend
attempts to acquire a Redis lock (\texttt{at\_risk:pipeline:lock},
\texttt{SET NX EX 600}).  If the lock is acquired, the pipeline runs as
an \texttt{asyncio.create\_task}; if not (another run is recent), the
existing cached rows are served immediately.  The TTL is deliberately
not deleted on completion, creating a natural 10-minute cooldown that
prevents redundant re-computation on each page refresh.

\subsubsection{Pipeline Steps}

Figure~\ref{fig:at_risk_pipeline} shows the full pipeline flow.

\begin{figure}[H]
  \centering
  \includegraphics[width=\linewidth]{fig11_at_risk_pipeline}
  \caption{At-risk detection pipeline: trigger, Redis lock, batch
    queries, LLM call, and DB upsert.}
  \label{fig:at_risk_pipeline}
\end{figure}

\begin{enumerate}
  \item \textbf{Identify at-risk students.}  One SQL query finds all
    students below the \texttt{AT\_RISK\_THRESHOLD} (default: 70\%).
    Only at-risk students are processed; the rest are ignored entirely.

  \item \textbf{Freshness gate.}  Students with
    \texttt{generated\_at} $< 600$~s and \texttt{ollama\_reachable=true}
    are skipped, avoiding redundant LLM calls for data that has not
    changed.

  \item \textbf{Stale record cleanup.}  Rows in
    \texttt{at\_risk\_explanations} for students who have recovered above
    the threshold are deleted.

  \item \textbf{Pre-flight Ollama check.}  One \texttt{GET /api/tags}
    call validates that phi3:mini is present before iterating students.

  \item \textbf{Per-student profile + LLM call.}  Three batch SQL
    queries build the student profile: (a)~attendance records joined with
    session and course data; (b)~average temperature and air quality on
    missed sessions, grouped by course; (c)~peer attendance rate per
    enrolled course.  The profile is formatted into a $<$500-token
    multi-course prompt.  One Ollama call returns a 3--4 sentence
    cross-course prose summary.  The result is upserted to
    \texttt{at\_risk\_explanations}.
\end{enumerate}

\subsubsection{Performance Optimisation (Phase~20)}

Table~\ref{tab:at_risk_perf} compares the original N+1 implementation
with the Phase~20 batch-query optimisation.

\begin{table}[H]
\centering
\caption{At-Risk Pipeline Performance: Before vs.\ After Phase~20}
\label{tab:at_risk_perf}
\begin{tabular}{@{}p{3.5cm}rr@{}}
\toprule
\textbf{Metric} & \textbf{Before} & \textbf{After} \\
\midrule
Students iterated     & All 35          & At-risk only $\sim$8 \\
Sensor queries/student & $N$ sessions $\times$ 2 & 1 batch JOIN \\
Peer-rate queries/student & 1 per course  & 1 GROUP BY \\
Ollama calls/student  & courses $+$ 1 ($\sim$7) & 1 \\
HTTP connections      & New TCP per call & Shared AsyncClient \\
Runtime (8 students)  & 14--16~min      & $\sim$2--3~min \\
\bottomrule
\end{tabular}
\end{table}

\subsubsection{Prompt Constraints}

The phi3-mini prompt is designed with strict constraints to prevent
harmful or inaccurate output: under 500 tokens per prompt; references
only provided data; never mentions health, personal life, family, or
psychology; no evaluative judgments; plain prose under 100 words.

\subsubsection{Frontend Auto-Poll}

The at-risk page polls \texttt{GET /api/at-risk} every 8~seconds while
any student has \texttt{summary\_explanation=null} (explanation not yet
generated by the background pipeline), then stops when all rows are
populated.  If \texttt{ollama\_reachable=false}, an amber warning card
is shown instead of the explanation.

Figure~\ref{fig:atrisk_screenshot} shows the at-risk page in operation.

\begin{figure}[H]
  \centering
  \includegraphics[width=\linewidth]{screenshot_atrisk}
  \caption{At-Risk page: student card list (left) and LLM explanation
    detail panel (right).}
  \label{fig:atrisk_screenshot}
\end{figure}

\subsection{Attendance Forecasting Pipeline}
\label{subsec:forecast_pipeline}

\subsubsection{Trigger and Classification}

The forecasting pipeline is triggered by \texttt{GET /api/forecasting}
with a Redis lock TTL of 1800~seconds (30~minutes).  Courses with fewer
than three ended sessions receive a marker row (\texttt{trend\_data=[]},
\texttt{generated\_at=now}) to prevent the frontend polling condition
from becoming unsatisfiable.

Figure~\ref{fig:forecast_pipeline} shows the pipeline flow.

\begin{figure}[H]
  \centering
  \includegraphics[width=\linewidth]{fig12_forecast_pipeline}
  \caption{Attendance forecasting pipeline: trend data query,
    deterministic classification, and Ollama interpretation.}
  \label{fig:forecast_pipeline}
\end{figure}

\subsubsection{Deterministic Trend Classification}

The trend classification algorithm operates on a chronological list of
per-session attendance rates (fractions 0--1):
\begin{enumerate}
  \item Compute consecutive deltas between adjacent rates.
  \item Compute \texttt{mean\_delta}.
  \item If \texttt{mean\_delta} $< -0.02$: classify as
    \texttt{steady\_decline}; if the second half of deltas is worse
    than the first half by $> 0.01$: upgrade to
    \texttt{accelerating\_decline}.
  \item If \texttt{mean\_delta} $> +0.015$: classify as
    \texttt{recovering}.
  \item Otherwise: \texttt{stable}.
\end{enumerate}

Confidence is set to \texttt{high} ($\geq$6 sessions),
\texttt{medium} ($\geq$4), or \texttt{low} ($<$4).

\subsubsection{LLM Role: Prose Only}

The LLM receives a compact 2-line structured prompt and returns exactly
two lines: \texttt{EXPECTED\_NEXT: <int>} (projected attendance
percentage) and \texttt{INTERPRETATION: <sentence>} (prose summary).
Classification and confidence labels are always computed
deterministically --- the LLM never produces the structured fields that
drive the dashboard's colour-coding logic, eliminating the risk of
hallucinated labels breaking the frontend.

Figure~\ref{fig:forecasting_screenshot} shows the forecasting page.

\begin{figure}[H]
  \centering
  \includegraphics[width=\linewidth]{screenshot_forecasting}
  \caption{Forecasting page: Recharts AreaChart with 70\% reference
    line, trend classification badge, and LLM interpretation card.}
  \label{fig:forecasting_screenshot}
\end{figure}

\subsection{Insights Engine}
\label{subsec:insights}

Beyond the LLM pipelines, a pure SQL analytics engine provides
on-demand charts with no AI latency: weekly attendance trend,
attendance heatmap by day-of-week and hour-slot, attendance decay
analysis, classroom comfort score (from Redis sensor averages), AC
effectiveness (temperature delta during active AC periods), and a
temperature-versus-attendance scatter plot.  All queries are
role-filtered, so professors see only their assigned courses.

\subsection{Design Decision: Deterministic Classification vs.\ LLM}
\label{subsec:det_vs_llm}

A key design decision is to use the LLM only for prose generation and
never for structured classification.  LLMs are probabilistic; the same
prompt can produce ``declining'' on one call and ``steady\_decline'' on
another.  The frontend uses the classification string to select badge
colours (red for accelerating decline, amber for steady decline, green
for recovering).  A hallucinated label would silently produce wrong
colours.

Delta mathematics, by contrast, is deterministic: the same input
always produces the same output, and execution time is microseconds
rather than seconds.  This separation --- deterministic structure,
LLM prose --- is the architectural pattern that makes the forecasting
feature robust enough for production use.
